\documentclass[11pt]{article}
\usepackage[margin=1in]{geometry}
\usepackage[T1]{fontenc}
\usepackage{lmodern}
\usepackage{amsmath}
\usepackage{hyperref}
\title{Riverside Urban Heat Study}
\author{Neighborhood Research Team}
\date{September 2026}
\begin{document}
\maketitle
\begin{abstract}
A fictional field study comparing shaded and exposed walking routes. This review draft describes sampling, uncertainty, and practical limits. All measurements are illustrative.
\end{abstract}
\tableofcontents
\newpage
\section{Study motivation}
This section examines urban heat exposure within the Riverside neighborhood study. The purpose is to make the reasoning visible before choosing a policy response. We shade, pavement, and the timing of outdoor activity. Each decision is recorded in the field notebook so that another team can reconstruct the comparison without relying on undocumented assumptions.

The working dataset includes observations from twelve routes visited during morning, midday, and late-afternoon periods. Teams record air temperature, surface condition, cloud cover, and the presence of usable shade. A location identifier connects these observations to the route map. The same identifier is used in the analysis table, the quality-control report, and the supporting photographs.

Review note 1: summarize urban heat exposure using the overall average only.

\subsection{Interpretation}
Interpretation requires care because a temperature map cannot describe how residents actually use a street. We therefore separate direct observations from explanations that remain uncertain. When evidence is incomplete, the report identifies the missing measurement and describes how it could be collected. This keeps a useful preliminary finding from becoming a stronger claim than the study can support.

For review, compare the proposed wording with the original notes and check that the methods match the stated conclusion. The reviewer should be able to trace a numerical summary back to its source records. Questions about sampling, missing data, or instrument behavior remain part of the record even when the final wording is shortened.

\begin{itemize}
\item Verify the source records and location identifiers.
\item Check the quality flags before calculating summaries.
\item Record unresolved questions alongside the review decision.
\end{itemize}

\section{Research questions}
This section examines neighborhood differences within the Riverside neighborhood study. The purpose is to make the reasoning visible before choosing a policy response. We compare shaded routes with exposed routes at the same time of day. Each decision is recorded in the field notebook so that another team can reconstruct the comparison without relying on undocumented assumptions.

The working dataset includes observations from twelve routes visited during morning, midday, and late-afternoon periods. Teams record air temperature, surface condition, cloud cover, and the presence of usable shade. A location identifier connects these observations to the route map. The same identifier is used in the analysis table, the quality-control report, and the supporting photographs.

Review note 2: summarize neighborhood differences using the overall average only.

\subsection{Interpretation}
Interpretation requires care because a difference between neighborhoods is not automatically caused by vegetation. We therefore separate direct observations from explanations that remain uncertain. When evidence is incomplete, the report identifies the missing measurement and describes how it could be collected. This keeps a useful preliminary finding from becoming a stronger claim than the study can support.

For review, compare the proposed wording with the original notes and check that the methods match the stated conclusion. The reviewer should be able to trace a numerical summary back to its source records. Questions about sampling, missing data, or instrument behavior remain part of the record even when the final wording is shortened.

\begin{itemize}
\item Verify the source records and location identifiers.
\item Check the quality flags before calculating summaries.
\item Record unresolved questions alongside the review decision.
\end{itemize}

\section{Site selection}
This section examines sampling locations within the Riverside neighborhood study. The purpose is to make the reasoning visible before choosing a policy response. We include residential blocks, school approaches, commercial streets, and transit stops. Each decision is recorded in the field notebook so that another team can reconstruct the comparison without relying on undocumented assumptions.

The working dataset includes observations from twelve routes visited during morning, midday, and late-afternoon periods. Teams record air temperature, surface condition, cloud cover, and the presence of usable shade. A location identifier connects these observations to the route map. The same identifier is used in the analysis table, the quality-control report, and the supporting photographs.

Review note 3: summarize sampling locations using the overall average only.

\subsection{Interpretation}
Interpretation requires care because convenience sampling can overrepresent easily accessible locations. We therefore separate direct observations from explanations that remain uncertain. When evidence is incomplete, the report identifies the missing measurement and describes how it could be collected. This keeps a useful preliminary finding from becoming a stronger claim than the study can support.

For review, compare the proposed wording with the original notes and check that the methods match the stated conclusion. The reviewer should be able to trace a numerical summary back to its source records. Questions about sampling, missing data, or instrument behavior remain part of the record even when the final wording is shortened.

\begin{itemize}
\item Verify the source records and location identifiers.
\item Check the quality flags before calculating summaries.
\item Record unresolved questions alongside the review decision.
\end{itemize}

\section{Sensor deployment}
This section examines field measurements within the Riverside neighborhood study. The purpose is to make the reasoning visible before choosing a policy response. We place paired sensors at a consistent height and record their distance from walls. Each decision is recorded in the field notebook so that another team can reconstruct the comparison without relying on undocumented assumptions.

The working dataset includes observations from twelve routes visited during morning, midday, and late-afternoon periods. Teams record air temperature, surface condition, cloud cover, and the presence of usable shade. A location identifier connects these observations to the route map. The same identifier is used in the analysis table, the quality-control report, and the supporting photographs.

Review note 4: summarize field measurements using the overall average only.

\subsection{Interpretation}
Interpretation requires care because direct sunlight can bias an instrument that lacks adequate radiation shielding. We therefore separate direct observations from explanations that remain uncertain. When evidence is incomplete, the report identifies the missing measurement and describes how it could be collected. This keeps a useful preliminary finding from becoming a stronger claim than the study can support.

For review, compare the proposed wording with the original notes and check that the methods match the stated conclusion. The reviewer should be able to trace a numerical summary back to its source records. Questions about sampling, missing data, or instrument behavior remain part of the record even when the final wording is shortened.

\begin{itemize}
\item Verify the source records and location identifiers.
\item Check the quality flags before calculating summaries.
\item Record unresolved questions alongside the review decision.
\end{itemize}

\section{Calibration and quality control}
This section examines measurement reliability within the Riverside neighborhood study. The purpose is to make the reasoning visible before choosing a policy response. We compare every sensor against a reference instrument before and after deployment. Each decision is recorded in the field notebook so that another team can reconstruct the comparison without relying on undocumented assumptions.

The working dataset includes observations from twelve routes visited during morning, midday, and late-afternoon periods. Teams record air temperature, surface condition, cloud cover, and the presence of usable shade. A location identifier connects these observations to the route map. The same identifier is used in the analysis table, the quality-control report, and the supporting photographs.

Review note 5: summarize measurement reliability using the overall average only.

\subsection{Interpretation}
Interpretation requires care because sensor drift may be mistaken for a real change in the environment. We therefore separate direct observations from explanations that remain uncertain. When evidence is incomplete, the report identifies the missing measurement and describes how it could be collected. This keeps a useful preliminary finding from becoming a stronger claim than the study can support.

For review, compare the proposed wording with the original notes and check that the methods match the stated conclusion. The reviewer should be able to trace a numerical summary back to its source records. Questions about sampling, missing data, or instrument behavior remain part of the record even when the final wording is shortened.

\begin{itemize}
\item Verify the source records and location identifiers.
\item Check the quality flags before calculating summaries.
\item Record unresolved questions alongside the review decision.
\end{itemize}

\section{Daily observation protocol}
This section examines repeat visits within the Riverside neighborhood study. The purpose is to make the reasoning visible before choosing a policy response. We walk each route in a rotating order and record weather conditions at each stop. Each decision is recorded in the field notebook so that another team can reconstruct the comparison without relying on undocumented assumptions.

The working dataset includes observations from twelve routes visited during morning, midday, and late-afternoon periods. Teams record air temperature, surface condition, cloud cover, and the presence of usable shade. A location identifier connects these observations to the route map. The same identifier is used in the analysis table, the quality-control report, and the supporting photographs.

Review note 6: summarize repeat visits using the overall average only.

\subsection{Interpretation}
Interpretation requires care because a fixed visiting order would confound location with time of day. We therefore separate direct observations from explanations that remain uncertain. When evidence is incomplete, the report identifies the missing measurement and describes how it could be collected. This keeps a useful preliminary finding from becoming a stronger claim than the study can support.

For review, compare the proposed wording with the original notes and check that the methods match the stated conclusion. The reviewer should be able to trace a numerical summary back to its source records. Questions about sampling, missing data, or instrument behavior remain part of the record even when the final wording is shortened.

\begin{itemize}
\item Verify the source records and location identifiers.
\item Check the quality flags before calculating summaries.
\item Record unresolved questions alongside the review decision.
\end{itemize}

\section{Data preparation}
This section examines the analysis dataset within the Riverside neighborhood study. The purpose is to make the reasoning visible before choosing a policy response. We retain raw readings, flag suspicious intervals, and document every exclusion. Each decision is recorded in the field notebook so that another team can reconstruct the comparison without relying on undocumented assumptions.

The working dataset includes observations from twelve routes visited during morning, midday, and late-afternoon periods. Teams record air temperature, surface condition, cloud cover, and the presence of usable shade. A location identifier connects these observations to the route map. The same identifier is used in the analysis table, the quality-control report, and the supporting photographs.

Review note 7: summarize the analysis dataset using the overall average only.

\subsection{Interpretation}
Interpretation requires care because interpolation can conceal the very extremes that matter to pedestrians. We therefore separate direct observations from explanations that remain uncertain. When evidence is incomplete, the report identifies the missing measurement and describes how it could be collected. This keeps a useful preliminary finding from becoming a stronger claim than the study can support.

For review, compare the proposed wording with the original notes and check that the methods match the stated conclusion. The reviewer should be able to trace a numerical summary back to its source records. Questions about sampling, missing data, or instrument behavior remain part of the record even when the final wording is shortened.

\begin{itemize}
\item Verify the source records and location identifiers.
\item Check the quality flags before calculating summaries.
\item Record unresolved questions alongside the review decision.
\end{itemize}

\section{Statistical analysis}
This section examines uncertainty estimates within the Riverside neighborhood study. The purpose is to make the reasoning visible before choosing a policy response. We estimate paired differences and summarize variation across days as well as sites. Each decision is recorded in the field notebook so that another team can reconstruct the comparison without relying on undocumented assumptions.

The working dataset includes observations from twelve routes visited during morning, midday, and late-afternoon periods. Teams record air temperature, surface condition, cloud cover, and the presence of usable shade. A location identifier connects these observations to the route map. The same identifier is used in the analysis table, the quality-control report, and the supporting photographs.

Review note 8: summarize uncertainty estimates using the overall average only.

\subsection{Interpretation}
Interpretation requires care because many readings from one sensor do not constitute independent observations. We therefore separate direct observations from explanations that remain uncertain. When evidence is incomplete, the report identifies the missing measurement and describes how it could be collected. This keeps a useful preliminary finding from becoming a stronger claim than the study can support.

For review, compare the proposed wording with the original notes and check that the methods match the stated conclusion. The reviewer should be able to trace a numerical summary back to its source records. Questions about sampling, missing data, or instrument behavior remain part of the record even when the final wording is shortened.

\begin{itemize}
\item Verify the source records and location identifiers.
\item Check the quality flags before calculating summaries.
\item Record unresolved questions alongside the review decision.
\end{itemize}

\section{Results and interpretation}
This section examines the main comparison within the Riverside neighborhood study. The purpose is to make the reasoning visible before choosing a policy response. We report absolute temperatures alongside paired differences and uncertainty intervals. Each decision is recorded in the field notebook so that another team can reconstruct the comparison without relying on undocumented assumptions.

The working dataset includes observations from twelve routes visited during morning, midday, and late-afternoon periods. Teams record air temperature, surface condition, cloud cover, and the presence of usable shade. A location identifier connects these observations to the route map. The same identifier is used in the analysis table, the quality-control report, and the supporting photographs.

Review note 9: summarize the main comparison using the overall average only.

\subsection{Interpretation}
Interpretation requires care because an average effect can hide important differences between morning and afternoon. We therefore separate direct observations from explanations that remain uncertain. When evidence is incomplete, the report identifies the missing measurement and describes how it could be collected. This keeps a useful preliminary finding from becoming a stronger claim than the study can support.

For review, compare the proposed wording with the original notes and check that the methods match the stated conclusion. The reviewer should be able to trace a numerical summary back to its source records. Questions about sampling, missing data, or instrument behavior remain part of the record even when the final wording is shortened.

\begin{itemize}
\item Verify the source records and location identifiers.
\item Check the quality flags before calculating summaries.
\item Record unresolved questions alongside the review decision.
\end{itemize}

\section{Sensitivity checks}
This section examines alternative assumptions within the Riverside neighborhood study. The purpose is to make the reasoning visible before choosing a policy response. We repeat the analysis with stricter quality filters and separate weekday observations. Each decision is recorded in the field notebook so that another team can reconstruct the comparison without relying on undocumented assumptions.

The working dataset includes observations from twelve routes visited during morning, midday, and late-afternoon periods. Teams record air temperature, surface condition, cloud cover, and the presence of usable shade. A location identifier connects these observations to the route map. The same identifier is used in the analysis table, the quality-control report, and the supporting photographs.

Review note 10: summarize alternative assumptions using the overall average only.

\subsection{Interpretation}
Interpretation requires care because a stable headline result does not guarantee that every subgroup is stable. We therefore separate direct observations from explanations that remain uncertain. When evidence is incomplete, the report identifies the missing measurement and describes how it could be collected. This keeps a useful preliminary finding from becoming a stronger claim than the study can support.

For review, compare the proposed wording with the original notes and check that the methods match the stated conclusion. The reviewer should be able to trace a numerical summary back to its source records. Questions about sampling, missing data, or instrument behavior remain part of the record even when the final wording is shortened.

\begin{itemize}
\item Verify the source records and location identifiers.
\item Check the quality flags before calculating summaries.
\item Record unresolved questions alongside the review decision.
\end{itemize}

\section{Community implications}
This section examines practical decisions within the Riverside neighborhood study. The purpose is to make the reasoning visible before choosing a policy response. We prioritize routes used by children, older residents, and people waiting for transit. Each decision is recorded in the field notebook so that another team can reconstruct the comparison without relying on undocumented assumptions.

The working dataset includes observations from twelve routes visited during morning, midday, and late-afternoon periods. Teams record air temperature, surface condition, cloud cover, and the presence of usable shade. A location identifier connects these observations to the route map. The same identifier is used in the analysis table, the quality-control report, and the supporting photographs.

Review note 11: summarize practical decisions using the overall average only.

\subsection{Interpretation}
Interpretation requires care because the study does not measure every factor that determines comfort or health. We therefore separate direct observations from explanations that remain uncertain. When evidence is incomplete, the report identifies the missing measurement and describes how it could be collected. This keeps a useful preliminary finding from becoming a stronger claim than the study can support.

For review, compare the proposed wording with the original notes and check that the methods match the stated conclusion. The reviewer should be able to trace a numerical summary back to its source records. Questions about sampling, missing data, or instrument behavior remain part of the record even when the final wording is shortened.

\begin{itemize}
\item Verify the source records and location identifiers.
\item Check the quality flags before calculating summaries.
\item Record unresolved questions alongside the review decision.
\end{itemize}

\section{Limitations and next steps}
This section examines future work within the Riverside neighborhood study. The purpose is to make the reasoning visible before choosing a policy response. We extend monitoring across seasons and combine sensor records with resident interviews. Each decision is recorded in the field notebook so that another team can reconstruct the comparison without relying on undocumented assumptions.

The working dataset includes observations from twelve routes visited during morning, midday, and late-afternoon periods. Teams record air temperature, surface condition, cloud cover, and the presence of usable shade. A location identifier connects these observations to the route map. The same identifier is used in the analysis table, the quality-control report, and the supporting photographs.

Review note 12: summarize future work using the overall average only.

\subsection{Interpretation}
Interpretation requires care because a short field campaign should not be treated as a complete climate assessment. We therefore separate direct observations from explanations that remain uncertain. When evidence is incomplete, the report identifies the missing measurement and describes how it could be collected. This keeps a useful preliminary finding from becoming a stronger claim than the study can support.

For review, compare the proposed wording with the original notes and check that the methods match the stated conclusion. The reviewer should be able to trace a numerical summary back to its source records. Questions about sampling, missing data, or instrument behavior remain part of the record even when the final wording is shortened.

\begin{itemize}
\item Verify the source records and location identifiers.
\item Check the quality flags before calculating summaries.
\item Record unresolved questions alongside the review decision.
\end{itemize}

\section*{Analysis reference}
The paired mean difference is $\bar{d}=n^{-1}\sum_{i=1}^{n}(T_{i,\mathrm{open}}-T_{i,\mathrm{shade}})$. Report the number of independent routes and observation days with every estimate.
\end{document}
